\section{Microscopic relaxation and the effective coefficient}
\label{sec:microscopic}

The weak gradient constructed in Section~\ref{sec:gradient} gives a coarse
lower bound on the limiting energy. We now identify its exact coefficient.
A correction that vanishes in $L^2$ can still alter the limiting rescaled
Dirichlet energy. We describe these corrections using polynomials in the unscaled
heat-bath operator $H$. Their limiting slope inner products are encoded by
a deterministic measure $\omega$ on $[0,\infty)$. Its atom at zero will
give the effective coefficient.

Throughout this section the couplings are Gaussian. Limits of quenched
quantities use the joint spectral convention of Section~\ref{sec:edge-law}.
Assertions about countably many fixed polynomials and cylinder functions
may consequently be imposed simultaneously after subsequence extraction.
All polynomial degrees are fixed before $n$ tends to infinity. Write
\[
 \langle A,B\rangle_{w,n}
   =\mu\!\left[\sum_i w_iA_iB_i\right],
 \qquad \|A\|_{w,n}^2=\langle A,A\rangle_{w,n}
\]
for vectors of functions on the cube. For a polynomial $Q$ and a fixed
leading eigenvector $v_a$, define $A_a^Q=DQ(H)Y_a$. Its component
$(A_a^Q)_i$ is independent of $X_i$.

\subsection{The measure of microscopic energies}

Let $\nu_{a,n}$ be the spectral measure of $H$ on $Y_a$, and put
$\omega_{a,n}(d\lambda)=\lambda\,\nu_{a,n}(d\lambda)$. This is a finite
positive measure. Selfadjointness and the formula for the Dirichlet form give
\begin{align}
 \int QQ'\,d\omega_{a,n}
   &=\langle A_a^Q,A_a^{Q'}\rangle_{w,n},
       \label{eq:micro-finite-measure}\\
 \int\lambda^j\,d\omega_{a,n}
   &=\mu\!\left[\sum_iw_iD_iY_aD_iH^jY_a\right].
       \label{eq:micro-finite-moments}
\end{align}
The factor $\lambda$ is essential: $\omega_{a,n}$ records energies, rather
than the variances recorded by the spectral measure of $Y_a$.

To specify the limiting measure through an experiment with deterministic
spectral data, use only fixed microscopic moments as follows.
Take deterministic midpoint quantiles
$\lambda_{1,N}\ge\cdots\ge\lambda_{N,N}$ of $\rho_{\rm sc}$, and choose
$y\in\R^N$ with $\|y\|^2=N$ such that
\[
 \frac1N\sum_{a=1}^N y_a^2\delta_{\lambda_{a,N}}
 \ \Longrightarrow\ \frac{d\rho_{\rm sc}(\lambda)}{2-\lambda}.
\]
For example, take $y_a^2$ proportional to $(2-\lambda_{a,N})^{-1}$ and
normalize their sum to $N$.  Draw $U$ by Haar measure subject to
$Uy=\mathbf1$, form $W=U\Lambda U^{\mathsf T}$, and define $H,w_i,Y_1$
as above, always omitting the diagonal of $W$ from the molecular fields.
The prescription for its $j$th moment, for each fixed integer $j\ge0$, is
\begin{equation}\label{eq:source-3}
 \int\lambda^j\,d\omega(\lambda)
 =\lim_{N\to\infty}\mathbb E_{U\mid Uy=\mathbf1}
 \left[\left.\sum_{i=1}^N
 w_iD_iY_1D_iH^jY_1\right|_{X=\mathbf1}\right],
 \qquad c_*=\omega(\{0\}).
\end{equation}
Proposition~\ref{prop:micro-moment-measure} proves that these limits exist
and determine a finite positive measure; Theorem~\ref{thm:micro-relaxation}
proves that its atom at zero is strictly positive and is the coefficient in
\eqref{eq:source-2}.

\begin{proposition}[Microscopic moment measure]
\label{prop:micro-moment-measure}
There is a deterministic finite positive measure $\omega$ on $[0,\infty)$
whose moments are given by \eqref{eq:source-3}. For fixed leading indices
$a,a'$, polynomials $Q,Q'$, and a bounded continuous cylinder $\psi$,
\begin{equation}
 \mu\!\left[\psi(x)\sum_iw_i(A_a^Q)_i(A_{a'}^{Q'})_i\right]
 \longrightarrow
 \1_{\{a=a'\}}\Pi[\psi]\int QQ'\,d\omega.
 \label{eq:source-32}
\end{equation}
For every $\varepsilon>0$ one also has
\begin{equation}
 \lim_{q\downarrow0}\limsup_{n\to\infty}
 \mathbb P\!\left(
  \mu\!\left[\left\|P_q(wA_a^Q)-v_a\int Q\,d\omega\right\|^2\right]
      >\varepsilon\right)=0.
 \label{eq:source-33}
\end{equation}
Here $wA$ denotes entrywise multiplication. For some $c>0$,
$\int e^{c\sqrt\lambda}\,d\omega(\lambda)<\infty$. The measure is uniquely
determined by its moments, and polynomials are dense in
$L^2((1+\lambda)\omega+\delta_0)$.
\end{proposition}

The proof has two parts. A fixed-order Haar calculation gives the limiting
contractions and their factorization. A separate estimate on iterated
differences controls moment growth. This separation avoids applying a
diagram expansion at a degree that increases with $n$.

\subsubsection{Conditional Haar contractions}

In the biased representation of Lemma~\ref{lem:edge-biased-law}, multiply
row $i$ of $U$ by the sampled spin $X_i$. The spin is then $\mathbf1$, and,
conditionally on $y=U^tX$, the rotation is Haar subject to $Uy=\mathbf1$.
The spectral inputs are
\[
 \frac1n\sum_a\delta_{\lambda_a}\Longrightarrow\rho_{\rm sc},
 \qquad
 \frac1n\sum_a y_a^2\delta_{\lambda_a}
      \Longrightarrow\frac{\rho_{\rm sc}(d\lambda)}{2-\lambda},
\]
together with bounded $\|\Lambda\|$, $\lambda_a\to2$, and
$y_a/\sqrt n\to0$ for each fixed leading index. These inputs hold both for
the deterministic data in \eqref{eq:source-3} and for a biased Gibbs sample.
For the latter, the weighted bulk limit is Lemma~\ref{lem:edge-bulk-weight}.

The case $Q=Q'=1$ shows what the calculation must track. Since
$D_iY_a=v_{ia}$, the unrelaxed slope contraction is
\[
 \sum_i\sech^2(h_i)v_{ia}v_{ia'}
   =\frac1n\sum_i\sech^2(h_i)
                     (\sqrt n\,v_{ia})(\sqrt n\,v_{ia'}).
\]
Thus each summand carries two marked eigenvector entries and the common
normalization is $1/n$. Its limit can already be read from a finite number
of rows. Put $u=y/\sqrt n$ and
\[
 m_n=u^t\Lambda u\longrightarrow1,\qquad
 s_n^2=u^t\Lambda^2u-m_n^2\longrightarrow1.
\]
These values follow from
$\int\lambda\,\rho_{\rm sc}(d\lambda)/(2-\lambda)=1$ and
$\int\lambda^2\,\rho_{\rm sc}(d\lambda)/(2-\lambda)=2$.
The exact field decomposition is
\[
 U\Lambda y=m_n\mathbf1+\sqrt n\,U(\Lambda u-m_nu).
\]
Conditionally on $Uy=\mathbf1$, the second term comes from a vector of
length $s_n$ rotated uniformly in $\mathbf1^\perp$.
It therefore tends at each fixed row to a standard normal.
The inner products of a marked axis $e_a$ with $u$ and with
$\Lambda u-m_nu$ are, respectively,
$y_a/\sqrt n$ and $(\lambda_a-m_n)y_a/\sqrt n$, both tending to zero.
Consequently the joint fixed-row limit of the field and the marked
entries is $(1+Z,G_a,G_{a'},\ldots)$, with independent standard normals
$Z,G_a,G_{a'},\ldots$ for distinct marked axes. The omitted diagonal
$W_{ii}$ tends to zero and does not change this limit.

The same conditional Haar argument at two distinct rows gives independent
copies of these limits, with uniformly bounded fixed moments.
Row exchangeability then makes the displayed average converge in $L^2$
to $\1_{\{a=a'\}}m_0$, where
\begin{equation}
 m_0=\mathbb E[\sech^2(1+Z)]=\omega([0,\infty)).
 \label{eq:micro-zeroth-moment}
\end{equation}
This is the energy before microscopic correction; the coefficient in the
limiting dynamics is instead the atom $c_*=\omega(\{0\})$.
For higher moments, a difference of a field function introduces a factor
$W_{ij}$ joining two sites. The graph calculation below keeps these
matrix factors together with the marked row variables.

A \emph{matrix edge} below means an entry of $UMU^t$, where $M$ is a bounded
spectral function of $\Lambda$ or a projection onto one of finitely many
marked spectral axes. A site vertex may carry a continuous function, of
polynomial growth, of the field $(U\Lambda y)_i$ and of the entries
$\sqrt n\,v_{ia}$ at the marked axes. In particular it may carry any fixed
derivative of $\tanh$.

Here are precise definitions for the contraction statement. Set
\[
 B=\operatorname{span}\{y/\sqrt n,\Lambda y/\sqrt n,
                                e_a:\ a\text{ is marked}\},
 \qquad R=\Id-P_B.
\]
Choose an orthonormal basis $b_1,\ldots,b_k$ of $B$, with
$b_1=y/\sqrt n$, and put
$\xi_i=(\sqrt n\,(Ub_r)_i)_{1\le r\le k}$.
The allowed vertex factors are functions of $\xi_i$; the field and
marked entries above are linear combinations with convergent coefficients.
If a finite multigraph $G=(V_G,E_G)$ has a symmetric matrix $M_e$ at each
edge and a vertex function $\Phi_v$, its normalized contraction is
\begin{equation}
 S_G=\frac1n\sum_{\kappa:V_G\to\{1,\ldots,n\}}
       \prod_{e=\{u,v\}\in E_G}
          (UM_eU^t)_{\kappa(u),\kappa(v)}
       \prod_{v\in V_G}\Phi_v(\xi_{\kappa(v)}).
 \label{eq:micro-graph-sum}
\end{equation}
Loops and multiple edges are allowed. We also use the same expression
with $\kappa$ restricted to be injective. By a normalized spectral trace
or a finite spectral path we mean, respectively,
\begin{equation}
 \frac1n\Tr(RM_{e_1}R\cdots M_{e_l}R),
 \qquad
 b_r^tM_{e_1}R\cdots RM_{e_l}b_s.
 \label{eq:micro-spectral-words}
\end{equation}
The one-edge path is $b_r^tM_e b_s$. In particular, path endpoints include
the field direction $\Lambda y/\sqrt n$, not just the spin direction.
Only bounded-length words occur in any fixed moment of $S_G$.

\begin{lemma}[Connected contractions]
\label{lem:micro-contractions}
Suppose the edge matrices have uniformly bounded operator norms and that
the traces and paths in \eqref{eq:micro-spectral-words} converge for every
fixed word. For a fixed connected $G$ and fixed continuous vertex
functions of polynomial growth, $S_G$ converges in every fixed $L^p$
under the conditional Haar law to a deterministic limit. The same holds
for the sum restricted to injective site labels. Products of finitely
many such contractions have the products of their limits.

The conclusion applies to the contractions defining
$\sum_iw_i(A_a^Q)_i(A_{a'}^{Q'})_i$ at the sampled spin, and to their
quadratic contractions with an additional $P_q$ edge, for $q>0$ fixed.
In the latter case the limit tends to zero as $q\downarrow0$ if that edge
is replaced by $P_q-P_{\rm mark}$, where $P_{\rm mark}$ projects onto the
marked spectral axes.
\end{lemma}

\begin{proof}
Separate all coincidences among the site labels; the $V$ remaining site
vertices have distinct labels. Condition on the image under $U$ of the
chosen orthonormal basis of $B$.
The residual map between the orthogonal complements is Haar. The Gram
matrix of the first two displayed vectors tends to
$\left(\begin{smallmatrix}1&1\\1&2\end{smallmatrix}\right)$, while the
marked axes are asymptotically orthonormal and orthogonal to these two
vectors. Thus the basis can be chosen with converging coefficients.
For bounds, an arbitrary orthonormal basis suffices.

At any fixed number of distinct rows, its image entries multiplied by
$\sqrt n$ converge jointly with all moments: the first column equals $1$,
and the remaining entries tend to independent standard normals.
This follows by representing finitely many columns by independent Gaussian
vectors projected off $\mathbf1$ and then orthonormalizing.
At fixed graph and moment orders, products of vertex factors and row
polynomials have a common polynomial growth bound. The bounds on higher
row moments make these products uniformly integrable, so their expectations
converge also when the continuous vertex functions are not polynomials.
The vertex functions can therefore be retained in the calculation of
residual Haar moments and integrated afterward.

We give the counting argument, using the orthogonal Haar integration
formula \cite[Section~3.2]{CollinsSniady2006}. Suppose that splitting the
edge endpoints between $B$ and $B^\perp$ leaves $e$ edges and $2b$ residual
endpoints. The finite-block endpoints contribute $n^{-(e-b)}$ times a
polynomial in the scaled row variables. Odd residual endpoint counts
have zero expectation. For $2b$ endpoints the integration formula is a
sum over a physical pairing $\pi$ and a spectral pairing $\sigma$, whose
coefficient has order
\[
 O(n^{-b-d(\pi,\sigma)}).
\]
Here $d(\pi,\sigma)$ is $b$ minus the number of loops in the paired union.
This order also follows directly by inverting the pairing Gram matrix:
after removing the diagonal power of $n$, the off-diagonal entries have
orders given by the pairing distances. The leading inverse terms are
paths of minimum total distance.

We spell out the spectral wiring represented by a pairing. Split each
matrix edge into its two endpoints, declaring each endpoint either
finite-block or residual. The $2b$ residual endpoints are paired by
$\sigma$; keep the original matrix edges as well. Each connected
component of this endpoint diagram is then either a closed alternating
cycle, or a path ending at two finite-block endpoints. Summing its
spectral indices gives a trace or path of
\eqref{eq:micro-spectral-words}. Define $t(\sigma)$ to be the number of
closed components of this diagram. The notation $t(\pi)$ means the
number obtained from exactly this construction with $\pi$ in place of
$\sigma$.

Writing $\mathsf P=\Id-UP_BU^t$, a term of the conditional expectation
therefore has the form
\[
 n^{-(e-b)}\,\mathrm{Wg}_{n-k}(\pi,\sigma)
 \prod_{\{\alpha,\beta\}\in\pi}
       \mathsf P_{\kappa(v_\alpha),\kappa(v_\beta)}
 \prod_{\text{cycles }c}\Tr(M_c)
 \prod_{\text{paths }p}\langle b_{r_p},M_p b_{s_p}\rangle,
\]
times its vertex factors and a polynomial in the scaled rows. Here
$v_\alpha$ is the physical vertex incident to endpoint $\alpha$, and
$M_c,M_p$ abbreviate the compressed products just defined.
There are $2e-2b$ finite-block endpoints, each supplying a row entry
$n^{-1/2}\xi_{i,r}$; this explains the factor $n^{-(e-b)}$.

The physical covariance $\mathsf P$ is the projection off the image of $B$.
Each pairing of unequal site vertices thus costs another factor $n^{-1}$,
with polynomially bounded factors in the scaled rows. Let $D$ be the
number of such pairs. Each trace costs at most $Cn$; each path costs at
most a constant. If the original graph has $a$
components, then
\begin{equation}
 t(\pi)\le e+D-V+a,
 \qquad
 t(\sigma)\le t(\pi)+d(\pi,\sigma).
 \label{eq:source-34}
\end{equation}
Indeed, add the $D$ unequal-site pairing connections to the graph, and
let $c\le a$ be the resulting number of components. Its cycle rank is
$e+D-V+c$. The residual trace circuits are edge-disjoint circuits in this
enlarged graph, so $t(\pi)\le e+D-V+c\le e+D-V+a$.
Open paths ending in $B$ add no circuits. A pairing switch creates at
most one new circuit, which proves the second inequality.

Before summing site labels, the resulting bound is
\[
 n^{-(e-b)}n^{-b-d(\pi,\sigma)}n^{-D}n^{t(\sigma)}
       =O(n^{-V+a}).
\]
It compensates for the site sum and the $a$ normalizing factors.
Every surviving term has a limit: the normalized trace and path factors
converge by hypothesis, the row expectations converge by Gaussian
orthogonalization, and the rescaled leading Haar coefficients converge
by the Gram-matrix expansion.

The same count proves concentration. In a product of normalized sums,
identifying vertices in two original components saves a power of $n$,
as does a physical pairing across components. If only the spectral
pairing crosses components, consider the first cross-component switch
along a shortest pairing path. Before it, each residual circuit or open
path lies inside one original component. Joining two circuits loses a
circuit; joining paths, or a path and a circuit, cannot create a circuit.
This switch therefore gains at most zero circuits, instead of the one
allowed by \eqref{eq:source-34}, and again saves a power of $n$.
Surviving terms remain inside the original components. Their row limits
and leading Haar coefficients factor; for the latter, shortest pairing
paths cannot join distinct paired components. Applying this factorization
to even powers of centered sums proves convergence in every fixed $L^p$.
The argument also gives bounds uniform over bounded spectral data.

We verify the graph representation for the heat-bath expressions.
Repeatedly use the exact rule
\[
 D_i(FG)=F(X)D_iG+(D_iF)G(X^{(i)}).
\]
A difference of an explicit spin identifies an existing site; a difference
of a field function adds a factor $W_{ij}$ joining the differentiating site
to that field vertex. The induction invariant is that each summand
obtained from one copy of $Y_a$ has a connected set of site labels and
exactly one marked factor. Initially its summand is $v_{ja}X_j$ on one
site. In an application of $H$, a newly summed site must differentiate
an existing factor: it either becomes an existing site or joins one by
a field edge. The factor $X_i-t_i$ remains at that same site. Flip
evaluations from the product rule introduce only field edges to sites
already present. The final difference $D_i$ has the same property, with
$i$ retained as the distinguished slope site. A term in which a new
difference strikes no factor vanishes, so there is no unattached
summation site. The slope inner product identifies the distinguished
sites of two such pieces; a projected quadratic contraction joins them
by a projection edge instead. Both are connected. Each piece carries
one marked factor $v_{ja}=n^{-1/2}(\sqrt n\,v_{ja})$, so the two marks
together supply the normalization $1/n$.

The shifts of fields caused by flips may be Taylor expanded to a
sufficiently large fixed order. Under conditional Haar, the off-diagonal
entries and centered diagonal entries are at most $n^{-0.49}$ outside a
probability exponentially small in a positive power of $n$. Haar
concentration on the complement of the prescribed spin direction has
bounded Lipschitz constants, and the entry means are $O(n^{-1})$;
a union bound gives the assertion. Each fixed expression has only
polynomially many summands. A large enough fixed Taylor order therefore
makes all remainders vanish in any prescribed fixed moment.
Polynomial bounds absorb the exceptional event. For the diagonal in
$h_i=(U\Lambda y)_i-W_{ii}$, Taylor-expand about $\Tr\Lambda/n$ and retain
the powers of $W_{ii}-\Tr\Lambda/n$ as loop factors in the graph.
Their edge matrix is $\Lambda-(\Tr\Lambda/n)\Id$; only the high-order
Taylor remainder is discarded by the entrywise bound.
This proves the required representation.

For these applications the edge matrices are powers of $\Lambda$,
fixed spectral cutoffs, marked projections, and the scalar shifts just
described. All hypotheses concerning \eqref{eq:micro-spectral-words}
follow from the stated spectral inputs. Indeed, finite-rank removal
changes a normalized trace by $O(n^{-1})$. Expanding each
$R=\Id-\sum_rb_rb_r^t$ in a path reduces it to finitely many products
of ordinary mixed spectral inner products between $y/\sqrt n$,
$\Lambda y/\sqrt n$, and the marked axes. Their limits are given by
the two empirical spectral laws, the leading eigenvalue limits, and
$y_a/\sqrt n\to0$. At fixed $q$, the cutoff endpoint $2-q^2$ is a
continuity point of the limiting laws.

Finally set $R_q=P_q-P_{\rm mark}$ in the extra edge. Normalized traces
containing $R_q$ tend to zero as $q\downarrow0$, because the semicircle
mass of $[2-q^2,2]$ vanishes. Finite paths also vanish: $R_q$ kills the
marked axes and, for every fixed $k$,
\[
 \limsup_{n\to\infty}\|R_q\Lambda^k y/\sqrt n\|^2
 \le \int_{2-q^2}^2\frac{\lambda^{2k}}{2-\lambda}
                                      \,\rho_{\rm sc}(d\lambda)
 \longrightarrow0.
\]
Paths containing finite-block compressions reduce to these statements by
expanding those projections. There are only finitely many contractions
at each fixed order. The limit order is $n\to\infty$ at fixed $q$, followed
by $q\downarrow0$.
\end{proof}

The lemma gives deterministic limits for polynomial slope inner products
at a sampled spin. A sign change of just one of two distinct marked axes
preserves the limiting spin/field Gram data and reverses the corresponding
contraction, so its limit is zero. All fixed leading axes have eigenvalue
limit $2$ and identical limiting Gram data; hence the diagonal limit is
independent of $a$. By \eqref{eq:micro-finite-measure}, it depends on $Q,Q'$
through $QQ'$.

The moment bounds in Lemma~\ref{lem:micro-contractions}, together with
Lemma~\ref{lem:edge-biased-law}, transfer convergence to constants from the
biased representation to quenched Gibbs integrals. Bounded cylinder
multipliers may then be inserted using
Corollary~\ref{cor:observations-static}. For \eqref{eq:source-33},
the lemma removes $P_q-P_{\rm mark}$. The remaining coefficients are
$v_{a'}\cdot(wA_a^Q)$, with diagonal limit given by the same moment
functional applied to $Q$ and off-diagonal limit zero.
It remains to prove that the moment functional determines a measure.

\subsection{Growth of the microscopic moments}
\label{subsec:micro-moment-growth}

Lemma~\ref{lem:micro-contractions} determines every fixed moment of the microscopic
spectral measure.  We next control their growth.  This will show that those
moments determine a measure on $[0,\infty)$ and that polynomials approximate
its atom at zero in the weighted norm needed for the correctors.  The
estimate is deterministic: it uses only a bound on the interaction norm and
small individual interaction entries.

For this subsection, let $W$ be a real symmetric $n\times n$ matrix with zero
diagonal, and let
\[
 \mu(X)=Z^{-1}\exp\bigl(\tfrac12 X^tWX\bigr),
 \qquad X\in\{-1,1\}^n.
\]
Write $h_i(X)=\sum_jW_{ij}X_j$, $t_i=\tanh h_i$, and $w_i=1-t_i^2$.
The half difference $D_i f$ is independent of $X_i$.  Thus the operator
$H=\sum_i(X_i-t_i)D_i$ satisfies
\begin{equation}
 \mu[fHg]=\sum_i\mu[w_iD_ifD_ig].
 \label{eq:micro-discrete-form}
\end{equation}
In particular, $H$ is selfadjoint and nonnegative on $L^2(\mu)$.
For an ordered tuple $S=(s_1,\ldots,s_m)$ of distinct sites, set
$D_S=D_{s_1}\cdots D_{s_m}$ and $w_S=\prod_{j\in S}w_j$, and define
\[
 N_m(f)^2=\sum_{\substack{S\in\{1,\ldots,n\}^m\\
                         \text{entries distinct}}}
                 \mu[w_S|D_Sf|^2].
\]
The order matters only in this sum: each set of $m$ sites occurs $m!$ times.
We put $N_m=0$ for $m>n$.  Subsets of $S$ in the proof below inherit the order
of their positions, whereas unions inside a difference operator denote the
corresponding collection of distinct sites.

\begin{lemma}[Weighted difference recursion]
\label{lem:micro-weighted-recursion}
Fix $K<\infty$.  Suppose that $W$ is symmetric, has zero diagonal, and obeys
\[
 \|W\|\le K,\qquad \epsilon_n:=\max_{i,j}|W_{ij}|\le n^{-49/100}.
\]
There is $C=C(K)\ge1$ such that, for each fixed integer $m\ge1$, uniformly
in these matrices and in real functions $f$ on $\{-1,1\}^n$,
\begin{equation}
 N_m(Hf)\le\sum_{l=1}^{m+2}
 \bigl(C^{m-l+3}(m+1)^{m-l+2}+o_{n;m}(1)\bigr)N_l(f).
 \label{eq:source-35}
\end{equation}
Here $o_{n;m}(1)$ is deterministic, nonnegative, and tends to zero as
$n\to\infty$ with $m$ fixed.
\end{lemma}

\begin{proof}
We first separate the differences that strike a conditional mean from those
that remain on $f$.  The exact product rule is
\[
 D_j(FG)=F D_jG+(D_jF)G(X^{(j)}),
 \qquad G(X^{(j)})=G(X)-2X_jD_jG(X).
\]
If $i\in S$, applying $D_i$ first to $(X_i-t_i)D_if$ gives $D_if$;
these terms therefore sum to $mD_Sf$.  For $i\notin S$, the terms in which
no difference strikes $t_i$ give $HD_Sf$.  Every remaining term, with its
overall minus sign, is
\begin{equation}
 (D_A t_i)\prod_{j\in B}(-2X_j)\,
 D_{\{i\}\cup(S\setminus A)\cup B}f,
 \qquad i\notin S,\quad \varnothing\ne A\subset S,\quad B\subset A.
 \label{eq:source-36}
\end{equation}
Indeed, the iterated product rule first evaluates
$D_{\{i\}\cup(S\setminus A)}f$ at $X^A$, and expanding the flips at the
sites of $A$ gives the displayed sum over $B$.

Fix the positions of $A$ and $B$, and write $a=|A|$, $b=|B|$ and
$l=1+m-a+b$.  We estimate the sum over $i$ of
\eqref{eq:source-36} in the norm obtained by summing
$\mu[w_S|\cdot|^2]$ over $S$.  For positive derivative orders,
\[
 |\tanh^{(a)}u|\le C^a a!\sech^2u.
\]
For example, Cauchy's formula on a fixed circle of radius less than
$\pi/2$, applied to $\sech^2z$, proves this inequality: on that circle
$|\sech^2z|\le C\sech^2u$ uniformly for real $u$.
The difference integral for $D_A t_i$ consequently gives
\begin{equation}
 |D_A t_i|\le w_i C^a a!\prod_{j\in A}|W_{ij}|.
 \label{eq:micro-coefficient-bound}
\end{equation}
The constant is uniform for each fixed $m$ and all sufficiently large $n$,
since all field shifts in that integral have absolute value at most
$2m\epsilon_n$ and $|\log\sech^2u-\log\sech^2v|\le2|u-v|$.

Suppose first that $b\ge1$, and select $j_0\in B$.  Apply Cauchy--Schwarz
to the sum over $i$, using $\sum_iW_{ij_0}^2\le K^2$ for its first factor.
The second factor contains $w_i^2$, the squared derivative of $f$, and
$\prod_{j\in A\setminus\{j_0\}}W_{ij}^2$.  Put $R=S\setminus A$ and
$A_0=A\setminus B$.  The weights $w_Rw_B$ occur in both input and output
norms and are retained.  Bound the additional output weights $w_{A_0}$
by one.  Summing the sites of $A_0$ now costs at most $K^{2(a-b)}$, by
row square sums; the remaining $b-1$ factors at sites of $B$ cost at most
$\epsilon_n^{2(b-1)}$.  Finally, $w_i^2\le w_i$ supplies the input weight
at site $i$.  Here only the distinctness restrictions involving the
summed-out sites of $A_0$ are removed; the retained tuple $(i,R,B)$
still consists of distinct sites.  After summing that tuple and integrating,
the resulting norm is at most
\begin{equation}
 C^{a+1}a!\epsilon_n^{b-1}N_l(f),\qquad b\ge1.
 \label{eq:micro-shared-slot-bound}
\end{equation}
Thus all terms with $b\ge2$ have coefficients $o_{n;m}(1)$.

When $b=0$, taking absolute values in the sum over $i$ would lose the
operator norm bound on $W$.  Instead expand the difference integral once
more:
\begin{align}
 D_A t_i={}&\Bigl(\prod_{j\in A}W_{ij}\Bigr)
 \left\{\tanh^{(a)}(h_i)
 -\tanh^{(a+1)}(h_i)\sum_{j\in A}W_{ij}X_j+R_{i,A}\right\},
 \label{eq:micro-coefficient-expansion}\\
 |R_{i,A}|\le{}&C_m w_i\epsilon_n^2.\nonumber
\end{align}
The first correction has the indicated sign because the field in the
integral is $h_i-2\sum_{j\in A}s_jW_{ij}X_j$, with each $s_j$ uniform on
$[0,1]$.

Here is the operator estimate for the first two terms.  Regard
\[
 (T_a u)_{j_1,\ldots,j_a}
       =\sum_i\prod_{r=1}^a W_{ij_r}\,u_i
\]
as a map from $\ell^2$ of sites to $\ell^2$ of ordered site tuples,
initially allowing repetitions.  Its Gram matrix is
$(WW^t)^{\circ a}$, the entrywise $a$th power, so $\|T_a\|\le K^a$.
To see the norm bound directly, the entrywise product of two matrices is
the compression of their tensor product to the span of the vectors
$e_i\otimes e_i$; iterate this observation for $WW^t$.
In each first-correction term, one factor $W_{ij_r}$ is replaced by
$W_{ij_r}^2$.  The matrix $(W_{ij}^2)_{i,j}$ has row and column sums at most
$K^2$, hence operator norm at most $K^2$.  The same tensor-product bound
therefore controls this corrected map by $K^{a+1}$.  Multiplication by
$X_{j_r}$ on its output is an isometry.

Fix $R=S\setminus A$ and apply these operator estimates pointwise in $X$,
before integrating against $w_R\mu$.  Restrictions that the sites of $A$ be distinct
and avoid $R$ are coordinate projections on the output of these maps.
The restriction $i\notin R$ is a coordinate projection on the input, and
$i\notin A$ is automatic from $W_{ii}=0$.  The output weights $w_A$ may
again be discarded.  The common weight $w_R$ remains, and the factors
$\tanh^{(a)}(h_i)$ and $\tanh^{(a+1)}(h_i)$ are bounded multiples of $w_i$.
Their multiplication on an input with squared norm
$\sum_iw_i|D_{\{i\}\cup R}f|^2$ thus costs no inverse powers of $w_i$.
Summing $R$ proves the required bound $C^{a+1}a!N_l(f)$ for the first
two terms of \eqref{eq:micro-coefficient-expansion}, after increasing $C$.
For the remainder, signs may be discarded: the matrix with entries
$\prod_{r=1}^a|W_{ij_r}|$ has squared Frobenius norm
\[
 \sum_i\left(\sum_jW_{ij}^2\right)^a\le nK^{2a}.
\]
Its contribution is at most
$C_m\sqrt n\epsilon_n^2N_l(f)=o_{n;m}(1)N_l(f)$.  We have proved
\begin{equation}
 \text{norm of the fixed-position terms with }b=0
 \ \le\bigl(C^{a+1}a!+o_{n;m}(1)\bigr)N_{1+m-a}(f).
 \label{eq:micro-unshared-slot-bound}
\end{equation}

It remains to control $HD_Sf$.  This term is measured with $w_S\mu$,
so we require an operator estimate in precisely that measure.  Fix $S$
and write $\nu=w_S\mu$, without normalizing it.  Denote its conditional
means and variances by $t_i'$ and $w_i'$, its heat-bath operator by
$H'=\sum_i(X_i-t_i')D_i$, and its first two difference seminorms by
$N_1'$ and $N_2'$, using the measure $\nu$ and weights $w_i'$.
We claim that
\begin{equation}
 \|HF\|_{L^2(\nu)}
 \le C\{N_2'(F)+(1+m)N_1'(F)\},
 \qquad C^{-1}w_i\le w_i'\le Cw_i.
 \label{eq:source-37}
\end{equation}
The constants depend only on $K$, once $n$ is sufficiently large for the
fixed $m$.

The conditional field under $\nu$ is exactly
\[
 h_i'=h_i+\delta_i,\qquad \delta_i=D_i\log w_S.
\]
Since the first two derivatives of $\log\sech^2u$ are bounded, the
nonnegative deterministic quantities
\[
 a_i=2\sum_{k\in S}|W_{ik}|,
 \qquad L_{ij}=C\sum_{k\in S}|W_{ik}W_{jk}|
\]
satisfy $|\delta_i|\le a_i$ and $|D_j\delta_i|\le L_{ij}$ for $i\ne j$.
They obey
\[
 \|a\|_2\le 2mK,\quad \|a\|_\infty\le2m\epsilon_n,
 \quad \|L\|\le CmK^2,
 \quad \Bigl(\sum_jL_{ij}^2\Bigr)^{1/2}\le CK a_i.
\]
These statements follow respectively from the triangle inequality for
columns of $W$ and the expression of $L$ as a sum of nonnegative
rank-one matrices.  The variance comparison in \eqref{eq:source-37}
follows from $|\delta_i|\le a_i=o(1)$ uniformly in $i$.  Taylor's formula
for $\tanh$ and $\sech^2$, now at the tilted fields, also gives
\begin{align}
 D_i t_j'&=w_j'(W_{ij}+E_{ij}),
 &|E_{ij}|&\le C(L_{ij}+W_{ij}^2),\label{eq:micro-tilted-coefficients}\\
 |D_iw_j'|&\le Cw_j' r_{ij},
 &r_{ij}&=|W_{ij}|+L_{ij},\qquad i\ne j.\nonumber
\end{align}
In deriving the first bound, the square of $L_{ij}$ is absorbed into
$L_{ij}$ since its maximum tends to zero.  The absolute error matrix is
bounded entrywise by a nonnegative matrix of norm at most $C(1+m)$;
indeed $(W_{ij}^2)$ has norm at most $K^2$.  Also
$\sup_i\sum_jr_{ij}^2\le C$ and $\max_{i,j}r_{ij}=o(1)$, for $n$ large
at fixed $m$.

For completeness, these coefficient bounds give the operator estimate
without an independence assumption on the slopes.  Put $F_i=D_iF$ and
$F_{ij}=D_iD_jF$.  The product rule gives the exact identity
\[
 D_iH'F=H'F_i+F_i
       -\sum_{j\ne i}(D_it_j')(F_j-2X_iF_{ij}).
\]
Use the version of \eqref{eq:micro-discrete-form} for $\nu$ to write
\begin{align*}
 \|H'F\|_{L^2(\nu)}^2
 ={}&N_1'(F)^2
 +\sum_{i\ne j}\nu[w_j'D_j(w_i'F_i)F_{ij}]\\
 &-\sum_{i\ne j}\nu[w_i'F_i(D_it_j')
                                      (F_j-2X_iF_{ij})].
\end{align*}
In the last sum, the $F_j$ term is bounded by
$C(1+m)N_1'(F)^2$: apply the signed norm bound to $W$ with input
$(w_i'F_i)_i$, and the absolute matrix bound to $(E_{ij})$.
The $F_{ij}$ term is bounded by $CN_1'(F)N_2'(F)$, by Cauchy--Schwarz
and the row square bounds for $r$.  Finally,
\[
 D_j(w_i'F_i)=w_i'F_{ij}+(D_jw_i')(F_i-2X_jF_{ij}).
\]
The first summand contributes $N_2'(F)^2$.  The other two summands cost,
respectively, $CN_1'(F)N_2'(F)$ and $CN_2'(F)^2$, using
\eqref{eq:micro-tilted-coefficients} and the same row square bounds.
Consequently
\[
 \|H'F\|_{L^2(\nu)}^2
 \le C\{N_2'(F)^2+N_1'(F)N_2'(F)+(1+m)N_1'(F)^2\}.
\]
Moreover $|t_i-t_i'|\le Ca_iw_i$, so pointwise Cauchy--Schwarz and
$\|a\|_2\le Cm$ show
\[
 \|(H-H')F\|_{L^2(\nu)}\le CmN_1'(F).
\]
These two bounds prove \eqref{eq:source-37}.

Apply \eqref{eq:source-37} to $F=D_Sf$ and sum squares over the ordered
tuples $S$.  A difference at a site already in $S$ vanishes.  For all
other sites, the variance comparison replaces $w_i'$ by $w_i$ with a
fixed factor.  Appending one or two ordered sites to $S$ enumerates
exactly the ordered tuples in $N_{m+1}$ or $N_{m+2}$; there is no further
factorial.  Minkowski's inequality therefore yields
\[
 \left(\sum_S\mu[w_S|HD_Sf|^2]\right)^{1/2}
       \le C\{N_{m+2}(f)+(1+m)N_{m+1}(f)\}.
\]
Together with the term $mD_Sf$, this controls every contribution except
the positional sums already bounded above.

There are $\binom ma$ choices of the positions of $A$, and
$\binom ab$ choices for $B$.  The terms with $b=0$ have
$l=m-a+1$ and coefficients bounded by $C^{a+1}a!\binom ma$;
those with $b=1$ have $l=m-a+2$ and coefficients bounded by
$C^{a+1}a!a\binom ma$.  Use $a!\binom ma\le m^a$ and absorb the
factor $a\le2^a$ by increasing $C$.  Both bounds, as well as the
$HD_Sf$ and $mD_Sf$ contributions, fit the coefficient
$C^{m-l+3}(m+1)^{m-l+2}$ in \eqref{eq:source-35}.  The finitely many
$b\ge2$ terms and Taylor remainders combine into $o_{n;m}(1)$.
\end{proof}

We now apply this estimate to the linear functions $Y_a=v_a^tX$,
where $\|v_a\|_2=1$.  Then $N_1(Y_a)\le1$ and $N_m(Y_a)=0$ for $m\ge2$.
For each fixed $j\ge0$, iterating \eqref{eq:source-35} gives
\begin{equation}
 \limsup_{n\to\infty}N_1(H^jY_a)
                  \le C^j(2j+2)^{2j}.
 \label{eq:micro-moment-growth}
\end{equation}
Here and below a larger constant $C=C(K)$ may be used.  To verify the
iteration explicitly, a nonzero terminal term follows indices
$m_0=1,m_1,\ldots,m_j=1$, with $1\le m_{r+1}\le m_r+2$.
Set $k_r=m_r-m_{r+1}+2\ge0$.  Then
\[
 \sum_{r=0}^{j-1}k_r=2j,\qquad
 \sum_{r=0}^{j-1}(k_r+1)=3j,\qquad m_r+1\le2j+2.
\]
The product of the coefficients on this path is at most
$C^{3j}(2j+2)^{2j}$.  The number of paths is bounded by the number of
weak compositions of $2j$ into $j$ parts,
$\binom{3j-1}{j-1}\le8^j$ for $j\ge1$.  Only finitely many fixed orders
occur for a given $j$, so all $o_{n;m}(1)$ coefficients disappear in this
iteration before $j$ is allowed to increase.  This proves
\eqref{eq:micro-moment-growth}; the case $j=0$ is immediate.

\begin{corollary}[Determination by moments and polynomial density]
\label{cor:micro-moment-density}
Let the matrices and normalized linear functions above vary with $n$,
and suppose the moments of the finite measures
\[
 \omega_{n,a}(B)=\langle Y_a,H\mathbf1_B(H)Y_a\rangle_{L^2(\mu)},
 \qquad B\subset[0,\infty)\text{ Borel},
\]
converge at every fixed nonnegative integer order.  They converge weakly
to a finite measure $\omega$ on $[0,\infty)$ with those moments.
This measure is uniquely determined among measures on $[0,\infty)$ by
its moments, and, for some $c>0$,
\[
 \int e^{c\sqrt\lambda}\,d\omega(\lambda)<\infty.
\]
Polynomials are dense in both $L^2(\omega)$ and
$L^2((1+\lambda)\omega+\delta_0)$.
\end{corollary}

\begin{proof}
The masses are at most one, and the spectral theorem and
\eqref{eq:micro-discrete-form} give
\[
 \int\lambda^{2j}\,d\omega_{n,a}(\lambda)=N_1(H^jY_a)^2.
\]
The $j=1$ bound gives tightness.  Bounds at higher even orders give
uniform integrability for each fixed power of $\lambda$, so every
weak subsequential limit has the prescribed moments.  In particular,
\eqref{eq:micro-moment-growth} implies
\begin{equation}
 \int\lambda^{2j}\,d\omega(\lambda)
                  \le C^{2j}(2j+2)^{4j}.
 \label{eq:micro-even-moments}
\end{equation}

Define the symmetric square-root lift $\rho$ by
\[
 \int g(t)\,d\rho(t)
   =\frac12\int\{g(\sqrt\lambda)+g(-\sqrt\lambda)\}\,d\omega(\lambda).
\]
Its $4j$th absolute moment is bounded by the right side of
\eqref{eq:micro-even-moments}.  For an arbitrary integer $k\ge1$, choose
$4j$ between $k$ and $k+3$ and use H\"older's inequality.  Since the mass
is at most one, this gives $\int|t|^k\,d\rho(t)\le(Dk)^k$ with a
constant $D=D(K)$.  The exponential series and $k!\ge(k/e)^k$ therefore
show that $\int e^{c|t|}\,d\rho(t)<\infty$ for small enough $c>0$.
This is the claimed square-root exponential moment.

If another measure on $[0,\infty)$ has the same moments, its symmetric
square-root lift has the same integer moments as $\rho$, and the same
argument gives an exponential moment.  The Fourier transforms of the
two lifts extend analytically to a strip about the real axis.  Their
Taylor series at zero agree, so the identity theorem and uniqueness of
Fourier transforms of finite measures identify the lifts, and hence the
original measures.  Every subsequential limit is therefore $\omega$,
proving weak convergence of the full sequence.

For density, let $g\in L^2(\rho)$ be orthogonal to all polynomials.
Cauchy--Schwarz implies that the finite signed measure $g\rho$ has an
exponential moment with a smaller positive exponent.  All derivatives
at zero of its Fourier transform vanish.  The same analytic argument
shows that this transform, and hence $g\rho$, is zero.  Thus polynomials
are dense in $L^2(\rho)$.  Applying the bounded projection
$p(t)\mapsto(p(t)+p(-t))/2$ shows that even polynomials are dense in its
even subspace.  The isometry $f(\lambda)\mapsto f(t^2)$ identifies this
subspace with $L^2(\omega)$, proving the first density assertion.
Finally, the symmetric square-root lift of
$(1+\lambda)\omega+\delta_0$ is $(1+t^2)\rho+\delta_0$.
It also has an exponential moment after decreasing $c$, and the same
proof establishes the second density assertion.
\end{proof}

\begin{proof}[Completion of the proof of Proposition~\ref{prop:micro-moment-measure}]
For the Gaussian matrices in this section, removing the diagonal preserves
a bounded operator norm and the entry bound used in
Lemma~\ref{lem:micro-weighted-recursion}, on events whose probability tends
to one. The fixed-order contractions give moment convergence there.
Corollary~\ref{cor:micro-moment-density} therefore constructs the unique
positive measure with those moments and proves the exponential bound and
polynomial density. The moments are deterministic and the same for every
fixed leading index. Applying the conditional Haar calculation to the
deterministic spectral data of \eqref{eq:source-3} gives precisely those
same moments. Finally the contraction and projection arguments preceding
this subsection give \eqref{eq:source-32} and \eqref{eq:source-33} with this
measure. This proves all assertions of the proposition.
\end{proof}



\subsection{Correctors and the exact energy}

Define $c_*=\omega(\{0\})$. We have proved that this number is deterministic;
for now we know only $c_*\ge0$. Choose real polynomials $Q_j$ with
\begin{equation}
 Q_j(0)=1,
 \qquad
 Q_j\longrightarrow\1_{\{0\}}
       \quad\hbox{in }L^2((1+\lambda)\omega+\delta_0).
 \label{eq:micro-zero-polynomials}
\end{equation}
The density assertion in Proposition~\ref{prop:micro-moment-measure}
supplies such polynomials: the extra atom enforces convergence of the
values at zero, after which division by those values normalizes them.
Consequently
\[
 \int Q_j^2\,d\omega\to c_*,
 \qquad \int Q_j\,d\omega\to c_*,
 \qquad \int\lambda Q_j^2\,d\omega\to0.
\]
The first limit gives the energy of the corrected slopes. The last makes
their microscopic divergence negligible.

\begin{lemma}[Cylinder recovery]
\label{lem:micro-recovery}
Let $F$ be a smooth compact cylinder, and let $Q$ be a polynomial with
$Q(0)=1$. Define
\begin{equation}
 f_{n,Q}(X)=F(x)+n^{-1/3}\sum_aF_{,a}(x)(Q(H)-1)Y_a,
 \label{eq:source-39}
\end{equation}
where the sum is over the finitely many coordinates of $F$.
Then $f_{n,Q}$ is bounded by a fixed power of $n$, converges weakly against
cylinders to $F$, and its squared $L^2$ norm converges to $\Pi[F^2]$.
Moreover,
\[
 \lim_{n\to\infty}T\cE(f_{n,Q})
       =\left(\int Q^2\,d\omega\right)\Pi[\|\nabla F\|^2].
\]
The exponent in the polynomial bound may depend on $Q$.
\end{lemma}

\begin{proof}
Set $V_a=(Q(H)-1)Y_a$. Since $Q-1$ vanishes at zero,
\[
 \|V_a\|_2^2
   =\int\frac{(Q(\lambda)-1)^2}{\lambda}
                                     \,d\omega_{a,n}(\lambda)
\]
is bounded for fixed $Q$: the quotient extends to a polynomial at zero.
Thus the correction in \eqref{eq:source-39} tends to zero in $L^2$.
Static convergence of cylinders gives the claimed convergence of values
and norms. The bounds $\|H\|_{\infty\to\infty}\le2n$ and
$\|Y_a\|_\infty\le\sqrt n$ give the polynomial supremum bound.

Taylor's formula in the finitely many coordinates of $F$ gives
\[
 n^{1/3}D_iF(x)=\sum_aF_{,a}(x)v_{ia}+r_i,
 \qquad \sum_i|r_i|^2=o(1),
\]
uniformly in the spin. The discrete product rule then writes the scaled
slope of \eqref{eq:source-39} as $\sum_aF_{,a}(x)A_a^Q$ plus this remainder
and terms $(D_iF_{,a})(x)V_a(X^{(i)})$. Their coefficients satisfy
\[
 |D_iF_{,a}(x)|^2\le Cn^{-2/3}\sum_bv_{ib}^2.
\]
The microscopic potentials need not be independent of these coefficients.
Conditionally on all spins except $i$, put
$p_i(\xi)=(1+\xi t_i)/2$. The identity
\[
 \frac{w_ip_i(\xi)}{p_i(-\xi)}=(1+\xi t_i)^2\le4
\]
implies
$\mu[w_i|V_a(X^{(i)})|^2]\le4\mu[|V_a(X)|^2]$.
Summing over $i$ and using the unit norms of the finitely many $v_b$ gives
\[
 \left\|n^{1/3}Df_{n,Q}-\sum_aF_{,a}(x)A_a^Q\right\|_{w,n}\to0.
\]
Equation~\eqref{eq:source-32}, with multiplier $F_{,a}F_{,a'}$, now proves
the energy limit.
\end{proof}

\begin{lemma}[Positivity of the effective coefficient]
\label{lem:micro-positivity}
One has $c_*=\omega(\{0\})>0$.
\end{lemma}

\begin{proof}
Choose a smooth compact cylinder $F$ with
$\Pi[\|\nabla F\|^2]>0$, possible because finite-coordinate marginals have
full support. For each fixed $j$, apply
Proposition~\ref{prop:gradient-observation} to $f_{n,Q_j}$.
Its limiting gradient is $\nabla F$ by
Proposition~\ref{prop:gradient-domain}, and the constant in the coarse
bound \eqref{eq:source-28} is independent of the fixed polynomial index.
Lemma~\ref{lem:micro-recovery} therefore yields
\[
 \Pi[\|\nabla F\|^2]
   \le C\left(\int Q_j^2\,d\omega\right)\Pi[\|\nabla F\|^2].
\]
Let $j\to\infty$ to obtain $1\le Cc_*$. Only one fixed-degree sequence is
used at a time. There is no need to impose a common polynomial supremum
exponent on a sequence whose degree increases with $n$.
\end{proof}

The remaining task is to test arbitrary finite-model slopes against these
corrected ones. Multiplication by a macroscopic cylinder test changes
their microscopic divergence by a negligible amount. Here $E_i$ denotes
conditional expectation over spin $i$, and $D_w^*$ is the adjoint of $D$
for the weighted slope pairing.

\begin{lemma}[Multiplied divergence]
\label{lem:micro-divergence}
For a fixed polynomial $Q$, fixed leading index $a$, and smooth compact
cylinder $\psi$,
\[
 D_w^*(\psi(x)A_a^Q)
       =\psi(x)HQ(H)Y_a+o_{L^2(\mu)}(1),
 \qquad
 \|HQ(H)Y_a\|_2^2\longrightarrow\int\lambda Q^2\,d\omega.
\]
\end{lemma}

\begin{proof}
Conditional single-site integration gives the exact formula
\[
 D_w^*(\psi A_a^Q)
     =\sum_i(X_i-t_i)(E_i\psi)(A_a^Q)_i.
\]
Replacing $E_i\psi$ by $\psi$ gives $\psi HQ(H)Y_a$.
The cylinder Lipschitz estimate is
$\sum_i|E_i\psi-\psi|^2\le Cn^{-2/3}$ pointwise. Since $(A_a^Q)_i$ is
independent of $X_i$, Cauchy--Schwarz gives
\[
 \begin{split}
 &\left\|\sum_i(X_i-t_i)(E_i\psi-\psi)(A_a^Q)_i\right\|_2^2\\
 &\quad\le Cn^{-2/3}
       \mu\!\left[\sum_i(X_i-t_i)^2|(A_a^Q)_i|^2\right]
       =Cn^{-2/3}\|A_a^Q\|_{w,n}^2=o(1).
 \end{split}
\]
Finally,
$\|HQ(H)Y_a\|_2^2=\int\lambda Q^2\,d\omega_{a,n}$.
\end{proof}

We can now state the form convergence used for the resolvent limit.
Weak cylinder convergence means convergence of the integrals against every
smooth compact cylinder test, in the quenched represented sense of
Section~\ref{sec:edge-law}. Strong convergence additionally requires
convergence of the squared $L^2$ norms.

\begin{theorem}[Relaxed energy]
\label{thm:micro-relaxation}
The constant $c_*=\omega(\{0\})$ is deterministic and strictly positive.
On any quenched represented subsequence, suppose that
$\|f_n\|_\infty\le Cn^D$ for fixed $C,D$, that the $L^2$ norms and
$T\cE(f_n)$ are bounded, and that $f_n$ converges weakly against cylinders
to $F$. Then $F$ is in the closed form domain on $L^2(\Pi)$ and
\begin{equation}
 \liminf_{n\to\infty}T\cE(f_n)
        \ge c_*\Pi[\|\nabla F\|_{\ell^2}^2].
 \label{eq:source-38}
\end{equation}
Conversely, for each smooth compact cylinder $F$, the functions
$f_{n,Q_j}$ in \eqref{eq:source-39} converge strongly to $F$ for every fixed
$j$, are polynomially bounded at that fixed $j$, and satisfy
\[
 \lim_{j\to\infty}\lim_{n\to\infty}T\cE(f_{n,Q_j})
                       =c_*\Pi[\|\nabla F\|^2].
\]
All assertions for countably many chosen sources and tests hold on a
common represented event. In particular the lower bound and recovery
use the same deterministic coefficient.
\end{theorem}

\begin{proof}
Only the lower bound remains. Pass to a subsequence attaining the lower
limit of the scaled energies. Proposition~\ref{prop:gradient-observation}
gives the weak gradient $G=\nabla F$ and its coarse bound. For finitely
many smooth compact cylinders $\psi_a$, put
$Z_{n,j}=\sum_a\psi_a(x)A_a^{Q_j}$. Equation~\eqref{eq:source-32} gives
\[
 \lim_{j\to\infty}\lim_{n\to\infty}\|Z_{n,j}\|_{w,n}^2
                         =c_*\sum_a\Pi[\psi_a^2].
\]
We identify its pairing with $n^{1/3}Df_n$.

Fix $q>0$ small and use the logarithmic averages over $[q,2q]$ in
Lemma~\ref{lem:gradient-smoothing}. By \eqref{eq:source-27},
$r_{n,q}=n^{1/3}(f_n-\overline{S f_n}_q)$ has $L^2$ norm at most $C/q$.
Lemma~\ref{lem:micro-divergence} bounds the limiting absolute value of its
slope pairing against $\psi_aA_a^{Q_j}$ by
\[
 \frac{C\|\psi_a\|_\infty}{q}
           \left(\int\lambda Q_j^2\,d\omega\right)^{1/2}.
\]
This vanishes when $j\to\infty$ at fixed $q$, after $n\to\infty$.

For the smoothed part, the second estimate of \eqref{eq:source-27}
replaces $n^{1/3}D\overline{S f_n}_q$ by $\overline{J^{f_n}}_q$,
with pairing error at most $Cq\limsup_n\|Z_{n,j}\|_{w,n}$.
This vector belongs to $\mathcal H_{2q}$ and has bounded unweighted
$L^2$ norm. For any fixed polynomial index $J$, \eqref{eq:source-33}
and weak convergence of the averaged slopes therefore give
\[
 \mu\!\left[\psi_a(x)\overline{J^{f_n}}_q\cdot wA_a^{Q_J}\right]
 \longrightarrow
       \left(\int Q_J\,d\omega\right)\Pi[\psi_aG_a]
\]
as $n\to\infty$ and then $q\downarrow0$ along the gradient extraction.

The last step fixes the polynomial before taking $q\downarrow0$, whereas
the residual estimate takes $j\to\infty$ first. The following bound
reconciles these orders without uniform graph estimates in the degree:
\[
 \limsup_{n\to\infty}
 \|w(A_a^{Q_j}-A_a^{Q_J})\|_{L^2(\mu;\R^n)}
 \le\left(\int(Q_j-Q_J)^2\,d\omega\right)^{1/2}.
\]
It follows from $0\le w_i\le1$ and \eqref{eq:source-32}.
For fixed $J$, its right side tends as $j\to\infty$ to
$\|\1_{\{0\}}-Q_J\|_{L^2(\omega)}$, and then tends to zero as
$J\to\infty$. Boundedness of $\overline{J^{f_n}}_q$ makes this estimate
uniform in the cap scale. All these bounds may be expressed with upper
and lower limits, so no prior convergence of the raw slope pairings is
required. We first take the finite-size limit, then send $j\to\infty$
at fixed $q$ to remove the residual, approximate the
remaining test by a large fixed $J$, send $q\downarrow0$, and finally send
$J\to\infty$. The limiting cross pairing is
$c_*\sum_a\Pi[\psi_aG_a]$.

Apply the elementary inequality
\[
 \|n^{1/3}Df_n\|_{w,n}^2
   \ge2\langle n^{1/3}Df_n,Z_{n,j}\rangle_{w,n}
                                      -\|Z_{n,j}\|_{w,n}^2
\]
and these iterated estimates to obtain
\[
 \liminf_nT\cE(f_n)
       \ge2c_*\sum_a\Pi[\psi_aG_a]-c_*\sum_a\Pi[\psi_a^2].
\]
Finite vectors of smooth compact cylinders are dense in
$L^2(\Pi;\ell^2)$. Optimizing proves \eqref{eq:source-38}.
The remaining claims follow from Lemmas~\ref{lem:micro-recovery}
and~\ref{lem:micro-positivity}.
\end{proof}

The coefficient in this theorem is the atom $\omega(\{0\})$, whereas
$\omega([0,\infty))$ is the unrelaxed energy of a leading linear coordinate.
Although $\omega_{a,n}(\{0\})=0$ for every finite system, its weak limit can
acquire mass from eigenvalues tending to zero. The correctors isolate
precisely this mass. Theorem~\ref{thm:micro-relaxation} and cylinder density
now supply the energy estimates needed to pass to Gaussian resolvents.
